\noindent Tables~\ref{tab:components-generic-derivative} and \ref{tab:components-generic-iderivative} record the action of \(\partial\) and \(i\partial\) on a generic complex paravector field \(Z\).
\begin{componenttable}[tab:components-generic-derivative]{\partial Z}
\componentrow{real scalar}{\partial_t a+\gradv\cdot\mathbf A}{}
\componentrow{imaginary scalar}{\partial_t b+\gradv\cdot\mathbf B}{}
\componentrow{real vector}{\gradv a+\partial_t\mathbf A-\gradv\times\mathbf B}{}
\componentrow{imaginary vector}{\gradv b+\gradv\times\mathbf A+\partial_t\mathbf B}{}
\end{componenttable}

\begin{componenttable}[tab:components-generic-iderivative]{i\,\partial Z}
\componentrow{real scalar}{-\partial_t b-\gradv\cdot\mathbf B}{}
\componentrow{imaginary scalar}{\partial_t a+\gradv\cdot\mathbf A}{}
\componentrow{real vector}{-\gradv b-\gradv\times\mathbf A-\partial_t\mathbf B}{}
\componentrow{imaginary vector}{\gradv a+\partial_t\mathbf A-\gradv\times\mathbf B}{}
\end{componenttable}

\noindent Tables~\ref{tab:components-fsquare} and \ref{tab:components-ffrev} summarize the electromagnetic invariants and the energy-density/Poynting split built from \(F\).
\begin{componenttable}[tab:components-fsquare]{\tfrac12F^{2}}
\componentrow{real scalar}{\tfrac12(\mathbf E^{2}-\mathbf B^{2})}{}
\componentrow{imaginary scalar}{i\,\mathbf E\cdot\mathbf B}{}
\componentrow{real vector}{\mathbf 0}{}
\componentrow{imaginary vector}{\mathbf 0}{}
\end{componenttable}

\begin{componenttable}[tab:components-ffrev]{\tfrac12F F^{\mathrm R}}
\componentrow{scalar}{\tfrac12(\mathbf E^{2}+\mathbf B^{2})}{energy density}
\componentrow{vector}{\mathbf E\times\mathbf B}{Poynting vector}
\end{componenttable}

\noindent Tables~\ref{tab:components-maxwell} and \ref{tab:components-wave-f} display the component splits of Maxwell's equation and its wave-equation consequence.
\begin{componenttable}[tab:components-maxwell]{\partial F=\rho-\mathbf J\cdot\sigv}
\componentrow{real scalar}{\gradv\cdot\mathbf E=\rho}{}
\componentrow{imaginary scalar}{\gradv\cdot\mathbf B=0}{}
\componentrow{real vector}{\partial_t\mathbf E-\gradv\times\mathbf B=-\mathbf J}{}
\componentrow{imaginary vector}{\partial_t\mathbf B+\gradv\times\mathbf E=0}{}
\end{componenttable}

\begin{componenttable}[tab:components-wave-f]{\square F=\partial^{*}(\rho-\mathbf J\cdot\sigv)}
\componentrow{real scalar}{\partial_t\rho+\gradv\cdot\mathbf J=0}{}
\componentrow{imaginary scalar}{0=0}{}
\componentrow{real vector}{\square\mathbf E=-\partial_t\mathbf J-\gradv\rho}{}
\componentrow{imaginary vector}{\square\mathbf B=\gradv\times\mathbf J}{}
\end{componenttable}

\noindent Table~\ref{tab:components-potential-gradient} records the component content of the potential identity \(\partial\Phi=S+F^{*}\).
\begin{componenttable}[tab:components-potential-gradient]{\partial\Phi=S+F^{*}}
\componentrow{real scalar}{\partial_tV+\gradv\cdot\mathbf A=S}{}
\componentrow{imaginary scalar}{0=0}{}
\componentrow{real vector}{\partial_t\mathbf A+\gradv V=-\mathbf E}{}
\componentrow{imaginary vector}{\gradv\times\mathbf A=\mathbf B}{}
\end{componenttable}

\noindent Table~\ref{tab:components-fu} is the component form of the product \(F U/\gamma\) used in the Lorentz-force law.
\begin{componenttable}[tab:components-fu]{F U/\gamma}
\componentrow{real scalar}{\mathbf E\cdot\mathbf v}{}
\componentrow{imaginary scalar}{\mathbf B\cdot\mathbf v}{}
\componentrow{real vector}{\mathbf E+\mathbf v\times\mathbf B}{}
\componentrow{imaginary vector}{\mathbf B+\mathbf E\times\mathbf v}{}
\end{componenttable}

\noindent Table~\ref{tab:components-partial-t} gives the direct component split of \(\partial T\) for the energy-density/Poynting paravector.
\begin{componenttable}[tab:components-partial-t]{\partial T}
\componentrow{real scalar}{\partial_t T_s+\gradv\cdot T_{\mathbf v}}{}
\componentrow{imaginary scalar}{0}{}
\componentrow{real vector}{\gradv T_s+\partial_t T_{\mathbf v}}{}
\componentrow{imaginary vector}{\gradv\times T_{\mathbf v}}{}
\end{componenttable}

\noindent Tables~\ref{tab:components-partial-fsquare}, \ref{tab:components-fdstar}, \ref{tab:components-dstarf}, and \ref{tab:components-dstarf-plus-fdstar} collect the remaining differentiated and source-coupled field combinations used in the energy-momentum discussion.
\begin{componenttable}[tab:components-partial-fsquare]{\partial F^{2}}
\componentrow{real scalar}{\partial_t(\mathbf E^{2}-\mathbf B^{2})}{}
\componentrow{imaginary scalar}{\partial_t\!\big(2\,\mathbf E\cdot\mathbf B\big)}{}
\componentrow{real vector}{\gradv(\mathbf E^{2}-\mathbf B^{2})}{}
\componentrow{imaginary vector}{\gradv\!\big(2\,\mathbf E\cdot\mathbf B\big)}{}
\end{componenttable}

\begin{componenttable}[tab:components-fdstar]{F D^{*}}
\componentrow{real scalar}{-\mathbf E\cdot\mathbf J}{}
\componentrow{imaginary scalar}{-\mathbf B\cdot\mathbf J}{}
\componentrow{real vector}{\rho\mathbf E+\mathbf B\times\mathbf J}{}
\componentrow{imaginary vector}{\rho\mathbf B-\mathbf E\times\mathbf J}{}
\end{componenttable}

\begin{componenttable}[tab:components-dstarf]{D^{*}F}
\componentrow{real scalar}{-\mathbf E\cdot\mathbf J}{}
\componentrow{imaginary scalar}{-\mathbf B\cdot\mathbf J}{}
\componentrow{real vector}{\rho\mathbf E-\mathbf B\times\mathbf J}{}
\componentrow{imaginary vector}{\rho\mathbf B+\mathbf E\times\mathbf J}{}
\end{componenttable}

\begin{componenttable}[tab:components-dstarf-plus-fdstar]{D^{*}F+F D^{*}}
\componentrow{real scalar}{-2\,\mathbf E\cdot\mathbf J}{}
\componentrow{imaginary scalar}{-2\,\mathbf B\cdot\mathbf J}{}
\componentrow{real vector}{2\rho\mathbf E}{}
\componentrow{imaginary vector}{2\rho\mathbf B}{}
\end{componenttable}
